\section{Топология Крулля на группе Галуа}

Напомним (3.9), что конечное расширение $\Om/F$ является расширением Галуа, если оно нормально и сепарабельно, т.\,е.\
каждый неприводимый многочлен $f\in F[X]$, имеющий корень в $\Om$, имеет в $\Om$ ровно $\deg f$ различных корней. Так
же (см.\ введение) определяется алгебраическое расширение Галуа. Например, $\Fsep$~--- расширение Галуа поля $F$.
Ясно, что $\Om/F$ является расширением Галуа тогда и только тогда, когда $\Om$~--- объединение конечных подрасширений
Галуа (каждый элемент лежит в поле разложения своего минимального многочлена).

\begin{proposition}[7.3]
Если $\Om/F$~--- расширение Галуа, то $\Om/M$~--- расширение Галуа для каждого промежуточного поля $F\subset M\subset\Om$.
\end{proposition}

\begin{proof}
Пусть $f\in M[X]$ неприводим и имеет корень $a\in\Om$. Минимальный многочлен $g$ элемента $a$ над $F$ раскладывается в
$\Om[X]$ на различные линейные множители. Так как $f$ делит $g$ в $M[X]$, то и $f$ раскладывается в $\Om[X]$ на
различные линейные множители.
\end{proof}

\begin{proposition}[7.4]
Пусть $\Om/F$~--- расширение Галуа и $F\subset E\subset\Om$. Тогда каждый $F$-гомоморфизм $E\to\Om$ продолжается до
$F$-изоморфизма $\Om\to\Om$.
\end{proposition}

\begin{proof}
Пусть $\varphi_E\colon E\to\Om$~--- $F$-гомоморфизм. Рассмотрим множество $\mathcal S$ пар $(M,\varphi_M)$, где
$E\subset M\subset\Om$~--- поле, а $\varphi_M\colon M\to\Om$~--- гомоморфизм, продолжающий $\varphi_E$; положим
$(M,\varphi_M)\le(M',\varphi_{M'})$, если $M\subset M'$ и $\varphi_{M'}|_M=\varphi_M$. Множество $\mathcal S$ непусто
(содержит $(E,\varphi_E)$), и у каждой цепи $(M_i,\varphi_i)$ есть верхняя грань $(\bigcup M_i,\bigcup\varphi_i)$
(объединение возрастающей цепи полей~--- поле, а согласованные гомоморфизмы склеиваются). По лемме Цорна есть
максимальный элемент $(K,\varphi_K)$.

\emph{$K=\Om$.} Пусть $\alpha\in\Om\setminus K$, $g$~--- минимальный многочлен $\alpha$ над $K$, $f$~--- над $F$. Так как
$g\mid f$ в $K[X]$, то $\varphi_K(g)\mid\varphi_K(f)=f$ ($\varphi_K$ применяется к коэффициентам и тождествен на $F$). Но
$f$ раскладывается в $\Om$ на линейные множители (нормальность), поэтому у $\varphi_K(g)$ есть корень $\beta\in\Om$, и
$\varphi_K$ продолжается до $K(\alpha)\to\Om$, $\alpha\mapsto\beta$ (2.1). Это противоречит максимальности.

\emph{Сюръективность.} Пусть $\varphi\colon\Om\to\Om$~--- полученный $F$-гомоморфизм, $a\in\Om$ и $f$~--- минимальный
многочлен $a$ над $F$. В $\Om$ ровно $\deg f$ корней $f$; $\varphi$ инъективен и переводит корни $f$ в корни $f$, значит,
$\varphi(\Om)$ тоже содержит $\deg f$ корней $f$, т.\,е.\ все, и в частности $a$. Итак, $\varphi$~--- изоморфизм.
\end{proof}

\begin{corollary}[7.5]
Пусть $\Om\supset E\supset F$ такие, как в 7.4. Если $E$ переходит в себя под действием всех элементов $\Aut(\Om/F)$, то
$E/F$~--- расширение Галуа.
\end{corollary}

\begin{proof}
Пусть $f\in F[X]$ неприводим с корнем $a\in E$, и $a_1,\dots,a_n$~--- его $n=\deg f$ различных корней в $\Om$. Есть
$F$-изоморфизм $F[a]\to F[a_i]\subset\Om$, $a\mapsto a_i$; по 7.4 он продолжается до $\sigma\in\Aut(\Om/F)$. Так как
$\sigma E=E$, то $a_i=\sigma(a)\in E$. Значит, $E/F$ нормально; сепарабельность наследуется от $\Om$.
\end{proof}

\paragraph{Подгруппы $G(S)$.} Пусть $\Om/F$~--- расширение Галуа и $G=\Aut(\Om/F)$. Для конечного подмножества
$S\subset\Om$ положим
\[
G(S)=\{\sigma\in G\mid \sigma s=s\ \text{для всех }s\in S\}=\bigcap_{s\in S}\Stab_G(s)=\Gal(\Om/F(S)).
\]
Будем говорить, что $S$ \emph{инвариантно} (или $G$-инвариантно), если $\sigma S=S$ для всех $\sigma\in G$.

\begin{lemma}[7.5a]
\begin{enumerate}[label=\textup{(\alph*)}]
\item Для конечного $S\subset\Om$ множество $\bar S=\bigcup_{\sigma\in G}\sigma S$ конечно и инвариантно; оно состоит из
всех корней в $\Om$ минимальных многочленов элементов $S$.
\item Если $S$ конечно и инвариантно, то $F(S)/F$~--- конечное расширение Галуа и $G(S)$~--- нормальная подгруппа $G$.
\item Каждое конечное подрасширение Галуа $E\subset\Om$ имеет вид $F(S)$ с конечным инвариантным $S$.
\end{enumerate}
\end{lemma}

\begin{proof}
(a) $\sigma(s)$~--- корень минимального многочлена $s$; обратно, каждый такой корень имеет вид $\sigma(s)$ по 7.4 (как в
доказательстве 7.5). Корней конечное число.
(b) $F(S)$ переходит в себя под действием $G$, поэтому оно Галуа по 7.5; конечно, так как порождено конечным числом
алгебраических элементов. Далее, $\tau G(S)\tau^{-1}=G(\tau S)=G(S)$: действительно, $\sigma$ оставляет на месте все точки
$\tau S$ $\iff$ $\tau^{-1}\sigma\tau$ оставляет на месте все точки $S$.
(c) $E=F(x_1,\dots,x_n)$; положим $S=\overline{\{x_1,\dots,x_n\}}$. По (a) $S$ инвариантно и состоит из корней минимальных
многочленов $x_i$; все они лежат в $E$, так как $E/F$ нормально. Значит, $E=F(S)$.
\end{proof}

\begin{proposition}[7.6]
На $G=\Aut(\Om/F)$ существует единственная структура топологической группы, для которой подгруппы $G(S)$ ($S\subset\Om$
конечно) образуют базу открытых окрестностей единицы. В этой топологии подгруппы $G(S)$ со инвариантным $S$ образуют базу
окрестностей единицы, состоящую из открытых нормальных подгрупп.
\end{proposition}

\begin{proof}
Проверим для семейства $\{G(S)\}$ условия (a)--(d) из 7.2. (a): $G(S_1)\cap G(S_2)=G(S_1\cup S_2)$. (b) и (c): каждое
$G(S)$~--- подгруппа, так что $G(S)G(S)=G(S)$ и $G(S)^{-1}=G(S)$. (d): по лемме 7.5a $G(\bar S)$ нормальна и
$G(\bar S)\subset G(S)$, так что $G(\bar S)=\sigma G(\bar S)\sigma^{-1}\subset\sigma G(S)\sigma^{-1}$. По 7.2 получаем
единственную топологию с нужными свойствами. Второе утверждение: всякая базисная окрестность $G(S)$ содержит $G(\bar S)$
со инвариантным $\bar S$, а такие $G(\bar S)$ нормальны.
\end{proof}

\begin{definition}
Топология из 7.6 называется \emph{топологией Крулля}. Группу $\Aut(\Om/F)$ с топологией Крулля обозначают $\Gal(\Om/F)$ и
называют \emph{группой Галуа} расширения $\Om/F$. Группа $\Gal(\Fsep/F)$ называется \emph{абсолютной группой Галуа}
поля $F$.
\end{definition}

Из 7.5a и 7.6 получаем удобные описания топологии Крулля:
\begin{itemize}[leftmargin=*]
\item \emph{Подгруппы $\Gal(\Om/E)$, где $E\subset\Om$~--- конечное подрасширение Галуа, образуют базу окрестностей
единицы из открытых нормальных подгрупп.}
\item Стабилизаторы $\Stab_G(\alpha)=\Gal(\Om/F(\alpha))$ открыты, а их сдвиги $\sigma\Stab_G(\alpha)$ образуют предбазу
топологии: базисные множества $\sigma G(S)=\bigcap_{s\in S}\sigma\Stab_G(s)$. Наглядно: \emph{$\sigma$ и $\tau$ близки,
если они одинаково действуют на большом конечном множестве элементов $\Om$.}
\end{itemize}

\begin{remark}[7.6a]
\begin{enumerate}[label=(\alph*),leftmargin=*]
\item Если $F$ совершенно, то $\Fsep$~--- алгебраическое замыкание $F$. Например, $\Q^{\rm sep}=\Qal$ и
$\F_p^{\rm sep}=\Fpb$ (поле $\F_p$ совершенно: многочлен $X^{p^n}-X$ сепарабелен).
\item Абсолютная группа Галуа определена лишь с точностью до внутреннего автоморфизма. Пусть $F'$~--- другое
сепарабельное замыкание и $\varphi_1,\varphi_2\colon F'\to\Fsep$~--- два $F$-изоморфизма. Они дают изоморфизмы
$\psi_i\colon\Gal(F'/F)\to\Gal(\Fsep/F)$, $\sigma\mapsto\varphi_i\sigma\varphi_i^{-1}$, и
$\psi_2(\sigma)=\varphi\,\psi_1(\sigma)\,\varphi^{-1}$, где $\varphi=\varphi_2\varphi_1^{-1}\in\Gal(\Fsep/F)$. Так как
$\varphi_1$ не единственен, канонического изоморфизма нет.
\end{enumerate}
\end{remark}

\begin{proposition}[7.7]
Пусть $\Om/F$~--- расширение Галуа и $E\subset\Om$~--- конечное подрасширение Галуа. Тогда ограничение
\[
\pi_E\colon\Gal(\Om/F)\to\Gal(E/F),\qquad\sigma\mapsto\sigma|_E,
\]
~--- непрерывный сюръективный гомоморфизм (на $\Gal(E/F)$~--- дискретная топология) с ядром $\Gal(\Om/E)$. В частности,
$\Gal(\Om/F)/\Gal(\Om/E)\cong\Gal(E/F)$.
\end{proposition}

\begin{proof}
Сюръективность: $\tau\in\Gal(E/F)$~--- это $F$-гомоморфизм $E\to\Om$, и по 7.4 он продолжается до $\sigma\in\Gal(\Om/F)$.
Ядро~--- $\Gal(\Om/E)=G(S)$, где $S$~--- конечное множество образующих $E$; оно открыто. Гомоморфизм в дискретную группу
с открытым ядром непрерывен: прообраз каждой точки~--- смежный класс ядра, т.\,е.\ открытое множество.
\end{proof}

\begin{lemma}[7.7a: проверка непрерывности]
Пусть $T$~--- топологическая группа и $\varphi\colon T\to\Gal(\Om/F)$~--- гомоморфизм. Следующие условия равносильны:
(i)~$\varphi$ непрерывен; (ii)~для каждого конечного подрасширения Галуа $E\subset\Om$ подгруппа
$\varphi^{-1}(\Gal(\Om/E))$ открыта; (iii)~для каждого такого $E$ композиция $\pi_E\circ\varphi\colon T\to\Gal(E/F)$
непрерывна. Достаточно проверять (ii) или (iii) для $E$ из любого семейства конечных подрасширений Галуа, в одном из
которых содержится каждое конечное подрасширение Галуа.
\end{lemma}

\begin{proof}
Для гомоморфизма топологических групп непрерывность равносильна непрерывности в единице (сдвиги~--- гомеоморфизмы), а
это в силу базы $\{\Gal(\Om/E)\}$ равносильно тому, что прообраз каждой $\Gal(\Om/E)$ содержит открытую окрестность
единицы, т.\,е.\ (7.1a(e)) открыт. (iii)$\iff$(ii): ядро $\pi_E\circ\varphi$ равно $\varphi^{-1}(\Gal(\Om/E))$, а
гомоморфизм в дискретную группу непрерывен тогда и только тогда, когда его ядро открыто. Наконец, при $E\subset E'$
подгруппа $\Gal(\Om/E')$ меньше, и открытость её прообраза влечёт открытость прообраза $\Gal(\Om/E)$ (7.1a(e)).
\end{proof}

\begin{theorem}[7.8]
Группа Галуа $G=\Gal(\Om/F)$ расширения Галуа $\Om/F$ компактна, хаусдорфова и вполне несвязна (т.\,е.\ её компоненты
связности~--- одноточечные множества).
\end{theorem}

\begin{proof}
\emph{1) $G$ хаусдорфова.} Пусть $\sigma\ne\tau$. Тогда $\sigma(\alpha)\ne\tau(\alpha)$ для некоторого $\alpha\in\Om$. Для
$S=\{\alpha\}$ открытые множества $\sigma G(S)$ и $\tau G(S)$ не пересекаются: элементы первого переводят $\alpha$ в
$\sigma(\alpha)$, второго~--- в $\tau(\alpha)$.

\emph{2) Вложение в произведение.} Пусть $S$ пробегает конечные инвариантные подмножества $\Om$. Подгруппа $G(S)$
нормальна и имеет конечный индекс: она~--- ядро действия $G$ на конечном множестве $S$, т.\,е.\ гомоморфизма
$G\to\Sym(S)$. Наделим конечные группы $G/G(S)$ дискретной топологией и рассмотрим
\[
\Phi\colon G\to P=\prod_S G/G(S),\qquad \sigma\mapsto(\sigma G(S))_S.
\]
Он инъективен: если $\sigma\in G(S)$ для всех $S$, то $\sigma$ тождествен на каждом $\bar{\{\alpha\}}$, т.\,е.\
$\sigma=\id$.

\emph{3) Топология Крулля индуцирована с $P$.} Базу топологии произведения дискретных пространств образуют
<<цилиндры>> $C=\{x\in P\mid x_{S_1}=\sigma_1G(S_1),\dots,x_{S_k}=\sigma_kG(S_k)\}$. Их прообразы:
\[
\Phi^{-1}(C)=\sigma_1G(S_1)\cap\dots\cap\sigma_kG(S_k),
\]
и это пересечение либо пусто, либо (если $\tau$ лежит в нём) равно $\tau\bigl(G(S_1)\cap\dots\cap G(S_k)\bigr)=\tau G(S)$,
$S=S_1\cup\dots\cup S_k$ (снова конечное инвариантное). Обратно, $\tau G(S)$~--- прообраз цилиндра с одной координатой.
Значит, прообразы базы $P$ образуют в точности базу $\{\tau G(S)\}$ топологии Крулля, т.\,е.\ $\Phi$~--- гомеоморфизм
$G$ на свой образ $\Phi(G)$ с топологией подпространства. (Вообще, если $\Phi\colon X\to Y$ инъективно, то семейство
$\{\Phi^{-1}(U)\mid U\text{ открыто в }Y\}$~--- топология на $X$, в которой $\Phi$~--- гомеоморфизм на образ, и прообразы
базы $Y$~--- её база: взятие прообраза сохраняет объединения и пересечения.)

\emph{4) Образ замкнут.} Для $S_1\subset S_2$ имеем $G(S_2)\subset G(S_1)$, и определено отображение
$\pi_{S_1S_2}\colon G/G(S_2)\to G/G(S_1)$. Положим
\[
E(S_1,S_2)=\{x\in P\mid \pi_{S_1S_2}(x_{S_2})=x_{S_1}\}.
\]
Это множество замкнуто: оно состоит из точек, где совпадают два непрерывных отображения $P\to G/G(S_1)$ (проекция на
$S_1$-ю координату и проекция на $S_2$-ю с последующим $\pi_{S_1S_2}$), а $G/G(S_1)$ хаусдорфово. Покажем, что
\[
\Phi(G)=\bigcap_{S_1\subset S_2}E(S_1,S_2),
\]
откуда $\Phi(G)$ замкнуто как пересечение замкнутых множеств.

Включение <<$\subset$>> очевидно: $\pi_{S_1S_2}(\sigma G(S_2))=\sigma G(S_1)$. Обратно, пусть $x=(x_S)$ лежит в
пересечении (такие $x$ назовём \emph{согласованными семействами}). Выберем представители: $x_S=\tau_SG(S)$.
\begin{itemize}[leftmargin=*]
\item \emph{Класс $x_S$ задаёт биекцию $\varphi_S=\tau_S|_S\colon S\to S$.} Она не зависит от выбора представителя: если
$\tau_SG(S)=\tau'_SG(S)$, то $\tau_S^{-1}\tau'_S\in G(S)$, т.\,е.\ $\tau_S$ и $\tau'_S$ совпадают на $S$. Так как $S$
инвариантно, а $\tau_S$ инъективен, $\varphi_S$~--- перестановка $S$.
\item \emph{Согласованность.} При $S_1\subset S_2$ условие $x\in E(S_1,S_2)$ означает $\tau_{S_2}\in\tau_{S_1}G(S_1)$,
т.\,е.\ $\varphi_{S_2}|_{S_1}=\varphi_{S_1}$.
\item \emph{Определение $\varphi$.} Для $\alpha\in\Om$ положим $\varphi(\alpha)=\varphi_S(\alpha)$, где $S$~--- любое
конечное инвариантное множество, содержащее $\alpha$ (например, $\bar{\{\alpha\}}$). Это корректно: если $\alpha\in S$ и
$\alpha\in S'$, то $S\cup S'$ тоже конечно и инвариантно, и по согласованности
$\varphi_S(\alpha)=\varphi_{S\cup S'}(\alpha)=\varphi_{S'}(\alpha)$.
\item \emph{$\varphi$~--- $F$-автоморфизм.} Для $a,b\in\Om$ возьмём конечное инвариантное $S$, содержащее $a,b,a+b,ab$. На
$S$ отображение $\varphi$ совпадает с автоморфизмом $\tau_S$, поэтому $\varphi(a+b)=\varphi(a)+\varphi(b)$ и
$\varphi(ab)=\varphi(a)\varphi(b)$. Для $c\in F$ множество $\{c\}$ инвариантно и $\varphi(c)=c$. Инъективность~--- как у
любого гомоморфизма полей; сюръективность: для $b\in\Om$ возьмём инвариантное $S\ni b$, тогда $b=\varphi_S(a)=\varphi(a)$
для некоторого $a\in S$, так как $\varphi_S$~--- перестановка $S$.
\item \emph{$\Phi(\varphi)=x$.} На каждом $S$ автоморфизмы $\varphi$ и $\tau_S$ совпадают, т.\,е.\
$\tau_S^{-1}\varphi\in G(S)$ и $\varphi G(S)=x_S$.
\end{itemize}

\emph{5) $G$ компактна.} $P$ компактно по теореме Тихонова, $\Phi(G)$ замкнуто в $P$, значит, компактно; $G$
гомеоморфна $\Phi(G)$.

\emph{6) $G$ вполне несвязна.} Каждая $G(S)$ открыта, а значит, и замкнута (7.1a). Пусть $U$~--- компонента связности
единицы. Тогда $U=(U\cap G(S))\sqcup(U\setminus G(S))$~--- разбиение на два открытых в $U$ множества; по связности
$U\subset G(S)$. Это верно для всех $S$, а $\bigcap_SG(S)=\{e\}$, поэтому $U=\{e\}$. Сдвигами получаем то же для любой
точки.
\end{proof}

\begin{remark}[7.8a]
\begin{enumerate}[label=(\alph*),leftmargin=*]
\item По 7.5a(b),(c) каждое конечное инвариантное $S$ порождает конечное подрасширение Галуа $E=F(S)$, и каждое такое
подрасширение имеет этот вид; при этом $G/G(S)\cong\Gal(E/F)$ по 7.7. Поэтому доказанное~--- это ровно утверждение введения: $\Phi$
из~\eqref{eq:Phi} вкладывает $\Aut(\Om/F)$ в $\prod_E\Gal(E/F)$ гомеоморфно на \emph{замкнутый} образ, состоящий из
согласованных семейств $(\sigma_E)$ ($\sigma_{E'}|_E=\sigma_E$ при $E\subset E'$). Это вид $G$ как обратного предела (7.26).
\item Ту же топологию можно получить, не используя подрасширения. Для $\alpha\in\Om$ пусть $R_\alpha$~--- конечное
множество корней минимального многочлена $\alpha$ в $\Om$. Тогда $G\subset\prod_{\alpha\in\Om}R_\alpha$,
$\sigma\mapsto(\sigma\alpha)_\alpha$, и образ~--- множество наборов $(s_\alpha)$ с $s_{a+b}=s_a+s_b$, $s_{ab}=s_as_b$,
$s_c=c$ ($c\in F$): такой набор задаёт $F$-гомоморфизм $\Om\to\Om$, а он сюръективен, так как переставляет конечное
множество корней каждого многочлена. Каждое из этих условий затрагивает не более трёх координат и потому задаёт
замкнутое множество. Базисные окрестности единицы здесь~--- $\{\sigma\mid\sigma\alpha_i=\alpha_i,\ i\le k\}=G(S)$, так что
топология та же. Иначе говоря, топология Крулля~--- это топология \emph{поточечной сходимости} действия на $\Om$ (с
дискретной топологией на $\Om$).
\end{enumerate}
\end{remark}

\begin{proposition}[7.9]
Для любого расширения Галуа $\Om/F$ имеем $\Om^{\Gal(\Om/F)}=F$.
\end{proposition}

\begin{proof}
Включение $F\subset\Om^{\Gal(\Om/F)}$ очевидно. Пусть $x\in\Om\setminus F$, и $E\ni x$~--- конечное подрасширение Галуа.
По основной теореме для конечных расширений $E^{\Gal(E/F)}=F$, так что некоторый $\tau\in\Gal(E/F)$ двигает $x$. По 7.7
$\tau=\sigma|_E$ для некоторого $\sigma\in\Gal(\Om/F)$, и $\sigma(x)\ne x$.
\end{proof}

Следующее утверждение~--- бесконечный вариант теоремы Артина 3.11.

\begin{proposition}[7.10]
Пусть $G$~--- группа автоморфизмов поля $E$ (т.\,е.\ подгруппа $\Aut(E)$), являющаяся компактной топологической
группой, и $F=E^G$. Если стабилизатор каждого элемента $E$ открыт в $G$, то $E/F$~--- расширение Галуа и
$G=\Gal(E/F)$, причём топология $G$ совпадает с топологией Крулля.
\end{proposition}

\begin{proof}
\emph{1) Конечные инвариантные множества.} Пусть $x_1,\dots,x_n\in E$ и $H_i=\Stab_G(x_i)$. Подгруппа $H_i$ открыта,
значит, имеет конечный индекс (7.1a(d)); так как $\sigma H_i\mapsto\sigma x_i$~--- биекция $G/H_i\to Gx_i$, орбита $Gx_i$
конечна. Положим $S=Gx_1\cup\dots\cup Gx_n$ (конечное $G$-инвариантное множество) и
\[
N=\bigcap_{s\in S}\Stab_G(s)=\bigcap_{\sigma\in G,\ i\le n}\sigma H_i\sigma^{-1}
\]
(так как $\Stab_G(\sigma x_i)=\sigma H_i\sigma^{-1}$). Это пересечение конечного числа открытых подгрупп, поэтому $N$
открыта и имеет конечный индекс; она нормальна, так как $\tau N\tau^{-1}=\bigcap_{s\in S}\Stab_G(\tau s)=N$ ($\tau$
переставляет $S$).

\emph{2) $E^N=F(S)$, и это конечное расширение Галуа с группой $G/N$.} Положим $K=F(S)$ и $M=E^N$; ясно, что $K\subset M$.
\begin{itemize}[leftmargin=*]
\item $G$ переводит $K$ в себя ($S$ инвариантно), и ядро действия $G$ на $K$ равно $N$ (элемент тождествен на $K$ тогда и
только тогда, когда он тождествен на $S$). Значит, конечная группа $G/N$ действует на $K$ точно, и
$K^{G/N}=K\cap E^G=F$. По теореме Артина 3.11 $K/F$~--- расширение Галуа с группой $G/N$ и $[K:F]=(G:N)$.
\item $G$ переводит $M$ в себя: если $m\in M$, $\sigma\in G$, $\nu\in N$, то
$\nu\sigma(m)=\sigma(\sigma^{-1}\nu\sigma)(m)=\sigma(m)$, так как $\sigma^{-1}\nu\sigma\in N$. Подгруппа $N$ действует на
$M$ тривиально, так что на $M$ действует конечная группа $\Gamma$~--- образ $G/N$ в $\Aut(M)$, причём $|\Gamma|\le(G:N)$ и
$M^\Gamma=M\cap E^G=F$. По теореме Артина $[M:F]=|\Gamma|\le(G:N)$.
\item Итак, $[K:F]=(G:N)\ge[M:F]$ и $K\subset M$, откуда $K=M$.
\end{itemize}

\emph{3) $E/F$~--- расширение Галуа.} Каждый элемент $x\in E$ лежит в конечном расширении Галуа $F(S)$ из шага 2 (с
$n=1$, $x_1=x$), поэтому его минимальный многочлен над $F$ сепарабелен и раскладывается в $F(S)\subset E$.

\emph{4) $G\subset\Gal(E/F)$, и вложение непрерывно.} Элементы $G$ тождественны на $F=E^G$, так что
$G\subset\Aut(E/F)=\Gal(E/F)$. По 7.7a достаточно проверить, что прообраз каждой $\Gal(E/F(T))$, $T\subset E$ конечно,
открыт; он равен $\bigcap_{t\in T}\Stab_G(t)$~--- конечному пересечению открытых подгрупп.

\emph{5) $G$ замкнута в $\Gal(E/F)$}: это компактное подмножество хаусдорфова пространства (7.8).

\emph{6) $G$ плотна в $\Gal(E/F)$.} Базу топологии образуют множества $\tau\Gal(E/M)$, $\tau\in\Gal(E/F)$,
$M\subset E$~--- конечное подрасширение Галуа. Пусть $M=F(x_1,\dots,x_n)$. Так как $M/F$ нормально, а $G\subset\Gal(E/F)$,
орбиты $Gx_i$ лежат в $M$; значит, для $S=\bigcup Gx_i$ имеем $M=F(S)$. По шагу 2 ограничение $G\to\Gal(M/F)$
сюръективно (его образ~--- группа порядка $(G:N)=[M:F]$). Значит, найдётся $\sigma\in G$ с $\sigma|_M=\tau|_M$, т.\,е.\
$\sigma\in\tau\Gal(E/M)$.

\emph{7)} Замкнутое и плотное подмножество совпадает со всем пространством: $G=\Gal(E/F)$. Тождественное отображение
$G\to\Gal(E/F)$~--- непрерывная биекция компакта на хаусдорфово пространство, т.\,е.\ гомеоморфизм.
\end{proof}

\begin{remark}[7.10a]
\begin{enumerate}[label=(\alph*),leftmargin=*]
\item Условие <<$G$~--- группа автоморфизмов>> означает, что $G$ действует на $E$ \emph{точно}. Если компактная группа
$G$ действует на $E$ через гомоморфизм $\rho\colon G\to\Aut(E)$ с ядром $K$ (и стабилизаторы открыты), то
$K=\bigcap_{x\in E}\Stab(x)$ замкнута, и предложение, применённое к $G/K$, даёт $\Gal(E/E^G)\cong G/K$, а не $G$.
\item Открытость стабилизаторов~--- это непрерывность действия $G\times E\to E$, если $E$ снабдить дискретной топологией.
\end{enumerate}
\end{remark}

\begin{aside}[7.11]
Не каждая компактная вполне несвязная группа является абсолютной группой Галуа некоторого поля. Например, конечная
абсолютная группа Галуа поля характеристики~$0$ имеет порядок $1$ или $2$. Точнее, по теореме Артина--Шрайера (1927),
если поле $F$ не алгебраически замкнуто, но имеет конечный индекс в своём алгебраическом замыкании $\bar F$, то $F$
вещественно замкнуто и $\bar F=F[\sqrt{-1}]$. (Ср.\ с 7.27: \emph{любая} такая группа~--- группа Галуа \emph{некоторого}
расширения Галуа.)
\end{aside}
